\DocumentMetadata{
  pdfversion=2.0,
  pdfstandard=ua-2,
  lang=en-US,
  tagging=on,
  tagging-setup={math/setup=mathml-SE,tabsorder=structure}
}
\documentclass[11pt,letterpaper]{article}
\usepackage[margin=0.68in,includefoot]{geometry}
\usepackage{newcomputermodern}
\usepackage{amsmath,amscd,mathtools}
\usepackage{tikz,adjustbox}
\providecommand{\square}{\mdlgwhtsquare}
\usepackage[hidelinks]{hyperref}
\hypersetup{
  pdftitle={Analysis PhD Exam, January 1989},
  pdfauthor={University of Florida Department of Mathematics},
  pdfsubject={Graduate mathematics examination},
  pdfdisplaydoctitle=true,
  bookmarksopen=true
}
\setcounter{secnumdepth}{0}
\setlength{\parindent}{0pt}
\setlength{\parskip}{0.28em}
\setlength{\emergencystretch}{3em}
\setlength{\leftmargini}{2.15em}
\setlength{\leftmarginii}{2.65em}
\setlength{\leftmarginiii}{3em}
\pagestyle{empty}
\raggedbottom
\allowdisplaybreaks
\NewTaggingSocketPlug{block/list/label}{examproblem}{%
  \tagstructbegin{tag=Lbl}\tagstructend%
  \tagstructbegin{tag=\LBody}%
  \tagstructbegin{tag=\ProblemHeadingTag,title-o={\ProblemHeadingText}}%
  \tagmcbegin{tag=\ProblemHeadingTag,actualtext={\ProblemHeadingText}}%
  #2%
  \tagmcend%
  \tagstructend%
}
\newenvironment{examproblems}[2]{%
  \def\ProblemHeadingTag{#2}%
  \AssignTaggingSocketPlug{block/list/label}{examproblem}%
  \begin{enumerate}%
  \setcounter{enumi}{#1}%
  \renewcommand{\labelenumi}{\textbf{\theenumi.}}%
  \setlength{\topsep}{0.35em}%
  \setlength{\partopsep}{0pt}%
  \setlength{\itemsep}{0.48em}%
  \setlength{\parsep}{0.18em}%
}{\end{enumerate}}
\newenvironment{examsubparts}{%
  \AssignTaggingSocketPlug{block/list/label}{default}%
  \begin{enumerate}%
  \setlength{\topsep}{0.18em}%
  \setlength{\partopsep}{0pt}%
  \setlength{\itemsep}{0.16em}%
  \setlength{\parsep}{0.1em}%
}{\end{enumerate}}
\newenvironment{examinstructions}{%
  \par\begingroup\itshape
}{%
  \par\endgroup\vspace{0.18em}
}
\makeatletter
\renewcommand\subsection{\@startsection{subsection}{2}{0pt}%
  {-1.25ex plus -.25ex minus -.1ex}%
  {0.55ex plus .1ex}%
  {\normalfont\large\bfseries}}
\renewcommand{\@maketitle}{%
  \newpage\null\vskip -1em%
  {\footnotesize\bfseries
    \href{https://www.ufl.edu/}{University of Florida}, \href{https://math.ufl.edu/}{Department of Mathematics}\par}%
  \vskip -0.5em%
  \tagstructbegin{tag=H1,title={Analysis PhD Exam, January 1989}}%
  \tagpdfparaOff%
  \tagmcbegin{tag=H1}%
    {\LARGE\bfseries\href{https://gma.math.ufl.edu/exams/analysis-phd/}{Analysis PhD Exam}, January 1989\par}%
  \tagmcend%
  \tagpdfparaOn%
  \tagstructend%
  \vskip 0.25em%
  \tagmcbegin{artifact}%
    \hrule height 1pt%
  \tagmcend%
  \vskip 1em%
}
\makeatother
\title{Analysis PhD Exam, January 1989}
\author{}
\date{}
\begin{document}
\maketitle
\thispagestyle{empty}
\begin{examinstructions}
Answer 6 out of the 7 problems. Justify all work. To obtain any partial credit, all work must be presented in a neat, logical and concise fashion.
\end{examinstructions}
\begin{examproblems}{0}{H2}
\def\ProblemHeadingText{Problem 1.}
\item
Let~\(f\) be a Borel-measurable function defined on \([0,1]\) with \(\infty>f(x)\ge 0\).
\begin{examsubparts}
\item[\textnormal{Part (a)}]
Prove that there is a non-increasing function \(f^*\) defined on \([0,1]\) such that
\[
m\{x:f^*(x)\ge\lambda\}=m\{x:f(x)\ge\lambda\}
\]
 for every \(\lambda\ge 0\), where~\(m\) is Lebesgue measure on \([0,1]\). (\(f^*\) is called a monotone rearrangement of~\(f\).)
\item[\textnormal{Part (b)}]
Show that \(\int f^*\,dm=\int f\,dm\).
\end{examsubparts}
\def\ProblemHeadingText{Problem 2.}
\item
Let \((X,\mathcal F,\mu)\) be a finite measure space, and let~\(f\) and~\(g\) be two measurable functions such that
\[
\mu(f\in A,\ g\in B)=\mu(f\in A)\mu(g\in B)
\]
 for every \(A,B\in\mathcal B(\mathbb R)\). Let \(\nu(A)=\mu(g\in A)\). Show that
\[
\int H(f,g)\,d\mu=\int\!\int H(f,t)\,\nu(dt)\,d\mu
\]
 for every positive \(\mathcal B(\mathbb R^2)\)-measurable function~\(H\).
\def\ProblemHeadingText{Problem 3.}
\item
Consider \(\ell^\infty=\{(a_n)_{n=1}^\infty:\sup_n|a_n|<\infty\}\). This space becomes a Banach space when equipped with the norm \(\|(a_n)_{n=1}^\infty\|=\sup\{|a_n|:n\ge 1\}\).
\begin{examsubparts}
\item[\textnormal{Part (a)}]
Show there is a linear functional \(T:\ell^\infty\to\mathbb R\) with the property: if \(\lim_{n\to\infty}a_n\) exists, then
\[
T[(a_n)_{n=1}^\infty]=\lim_{n\to\infty}a_n.
\]
 (Use the Hahn–Banach theorem.)
\item[\textnormal{Part (b)}]
Modify your proof in (a) to show there is a linear functional \(S:\ell^\infty\to\mathbb R\) with the property: if \(\lim_{n\to\infty}n^{-1}\sum_{k\le n}a_k\) exists, then
\[
S[(a_n)_{n=1}^\infty]=\lim_{n\to\infty}n^{-1}\sum_{k\le n}a_k.
\]
\end{examsubparts}
\def\ProblemHeadingText{Problem 4.}
\item
State and prove the Radon–Nikodym theorem.
\def\ProblemHeadingText{Problem 5.}
\item
State and prove Fubini’s theorem.
\def\ProblemHeadingText{Problem 6.}
\item
Let~\(\mu\) be a measure on \(((0,\infty),\mathcal B(0,\infty))\) defined by
\[
\mu(A)=\int_A |x|^{-1/2}\,dx.
\]
 Let \(T:(0,\infty)\to(0,\infty)\) be given by \(T(x)=x^{-1}\). Define a new measure~\(\pi\) on \((0,\infty)\) by setting \(\pi(A)=\mu(T^{-1}(A))\). Compute \(d\pi/d\mu\).
\def\ProblemHeadingText{Problem 7.}
\item
Let \(c(t)\) be a function defined on \([0,\infty)\) with the following properties: \(c(t)\) is strictly increasing and continuous, \(c(0)=0\), and \(c(\infty)=\infty\). Let \(\tau_t=\inf\{s:c(s)>t\}\). Prove that
\[
\int_0^\infty f(t)\,c(dt)=\int_0^\infty f(\tau_t)\,dt
\]
 for every positive continuous function~\(f\) on \([0,\infty)\).
\end{examproblems}
\end{document}
